\documentclass[11pt]{article}
\usepackage[T1]{fontenc}
\usepackage[utf8]{inputenc}
\usepackage{lmodern,amsmath,amssymb,microtype}
\usepackage[margin=0.83in]{geometry}
\usepackage{xcolor}
\usepackage[colorlinks=true,urlcolor=blue!55!black]{hyperref}
\setlength{\emergencystretch}{3em}
\setlength{\parindent}{0pt}
\setlength{\parskip}{0.30em}
\urlstyle{same}
\hypersetup{pdfauthor={Wentao Li},pdftitle={A finite generating list for shallow ReLU rank constraints}}
\title{A finite generating list for shallow ReLU rank constraints}
\author{Wentao Li\\{\small Prepared with assistance from OpenAI Codex}}
\date{November 20, 2026}
\begin{document}
\maketitle
\begin{abstract}
Integer simplex grids on rational subspaces generate the entire single-combination rank ideal. Linear algebra selects its contribution to the minimal pattern-ideal generators.
\end{abstract}
\section*{Which integer combinations suffice?}
For a one-hidden-layer network with $h$ hidden units, select $k$ activation patterns and write $A\in\{0,1\}^{h\times k}$ for the matrix with these patterns as columns. Let $M_1,\ldots,M_k$ be formal $q$-by-$d$ matrices of output linear coefficients. In the network parametrization each $M_i$ is the sum of the rank-one neuron matrices active in column $i$. For $\lambda\in\mathbb Z^k$ put
\[
 B(\lambda)=\sum_i\lambda_i M_i,\qquad s(\lambda)=|\mathrm{supp}(A\lambda)|.
\]
Theorem~34 of Alexandr and Mont\'ufar imposes the minors of size $s(\lambda)+1$ of $B(\lambda)$. Their Problem~35 asks for finitely many integer vectors supplying all minimal generators arising from these constraints. Write $I_{\rm rank}$ for the ideal generated by this family in the complex coordinate ring of the blocks and $J_A$ for the full homogeneous pattern ideal.

\paragraph{Theorem and construction.} Put $m=\min(d,q)$. For each $S\subseteq[h]$ with $|S|<m$, compute a rational basis of
\[
 V_S=\{v\in\mathbb Q^k:(Av)_j=0\ \text{for }j\notin S\}
\]
and clear denominators of its columns to obtain an integer matrix $U_S$ with $e=\dim V_S$ columns. For $r=|S|+1$ include the integer vectors
\[
 U_S t,\qquad t\in\mathbb Z_{\ge0}^{e},\quad \sum_j t_j\le r .
\]
The union over $S$ is a finite answer to Problem~35. Vectors equal to zero and duplicates may be dropped. The usual Theorem~34 minors at these vectors generate $I_{\rm rank}$ as an ideal, and their classes span all contributions of that family among minimal generators of $J_A$.

\section*{Interpolation on each rational subspace}
Fix $S$ and a choice of $r$ rows and columns of $B(U_S t)$. Its determinant $P(t)$ is a polynomial of total degree at most $r$ in $e$ variables with coefficients in the block-coordinate ring. Values on the integer simplex $|t|_1\le r$ determine such polynomials by linear interpolation over $\mathbb Q$. Indeed the products $\prod_j\binom{t_j}{\nu_j}$ for $|\nu|_1\le r$ form a basis. Their evaluation matrix on the same multi-indices, ordered by total degree, is triangular with ones on its diagonal. Interpolation therefore works with coefficients in the coordinate ring as well.

At a sample $t$ the actual support size $s=s(U_S t)$ may drop below $r-1$. The $r$-minor there belongs to the ideal generated by the $(s+1)$-minors at that same vector by repeated Laplace expansion. Thus every sampled value of $P$ belongs to the proposed finite ideal. For any integer $\lambda$ whose support under $A$ is $S$, write $\lambda=U_S t_0$ with $t_0$ rational. Interpolating $P(t_0)$ from the sampled values proves membership for all its minors. Larger support sizes produce no minors. Conversely every proposed generator is an original constraint. This proves the ideal equality, without radicalization.

\section*{Extracting the minimal contribution}
Let $\mathfrak m$ be the positive-degree coordinate ideal. Images in $J_A/\mathfrak m J_A$ measure independent minimal generators degree by degree. The ideal equality just proved shows that the image of the entire rank family is spanned by the finite sample minors there. Indeed, every expression in those generators reduces modulo $\mathfrak m J_A$ to a linear combination with constant coefficients. Selecting a basis of this span specifies which sampled generators to retain in a minimal homogeneous generating set of $J_A$.

This selection is effective: in each degree $j\le m$, substitute the given neuron parametrization on the space of degree-$j$ monomials in the block variables and take the kernel by finite coefficient linear algebra. This gives $(J_A)_j$. Multiplying $(J_A)_{j-1}$ by the coordinate variables gives $(\mathfrak m J_A)_j$. Row reduction then selects representatives from the sample minors. Using $I_{\rm rank}$ in place of $J_A$ similarly produces a minimal set for the rank ideal itself.\hfill$\square$

\paragraph{Size and scope.} A support of size $s$ contributes at most $\binom{e+s+1}{s+1}$ vectors. The construction covers nonzero kernels of $A$ and their linear equations. Whether the single-combination minors define the entire pattern variety is a separate matter. The source treats additional families involving concatenated matrices; the two-block ideal was subsequently determined by Alexandr, Grillo, Holleben, Mont\'ufar and Yu.

\section*{Sources}
\begin{itemize}
\item \href{https://arxiv.org/abs/2508.03867}{Yulia Alexandr and Guido Mont\'ufar (2025), Problem 35 and Theorem 34}.
\item \href{https://arxiv.org/abs/2611.20346}{Finkel, Fu and Wu, Canonical activation patterns and ReLU pattern varieties}.
\item \href{https://arxiv.org/abs/2611.23846}{Alexandr et al., Defining equations for two-block shallow ReLU pattern varieties}.
\end{itemize}
{\small Research draft. Independent verification of the proof and its priority is pending.}
\end{document}
